% Figure: physical-component schematic of the vapour-compression cycle:
% compressor, condenser, expansion valve, evaporator arranged as a loop.
\begin{figure}[htbp]
\centering
\begin{tikzpicture}[
  comp/.style={draw, very thick, regular polygon, regular polygon sides=3,
    shape border rotate=30, minimum size=1.3cm, fill=englime!15},
  hx/.style={draw, very thick, minimum width=2.6cm, minimum height=1.0cm, font=\small},
  valve/.style={draw, very thick, regular polygon, regular polygon sides=3,
    minimum size=0.9cm, fill=engorange!20},
  flow/.style={-{Latex[length=2.4mm]}, very thick, engblue},
]
% layout: condenser top, evaporator bottom, compressor right, valve left
\node[hx, fill=engred!12] (cond) at (0,2.4) {condenser};
\node[hx, fill=engblue!12] (evap) at (0,-2.4) {evaporator};
\node[comp] (comp) at (4,0) {};
\node[font=\scriptsize] at (4,0) {comp};
\node[valve] (valve) at (-4,0) {};
\node[font=\scriptsize, anchor=east] at (-4.6,0) {valve};
% loop arrows (clockwise): evap -> comp -> cond -> valve -> evap
\draw[flow] (evap.east) -| node[pos=0.25, below, font=\scriptsize]{1} (comp.south);
\draw[flow] (comp.north) |- node[pos=0.75, above, font=\scriptsize]{2} (cond.east);
\draw[flow] (cond.west) -| node[pos=0.75, above, font=\scriptsize]{3} (valve.north);
\draw[flow] (valve.south) |- node[pos=0.25, below, font=\scriptsize]{4} (evap.west);
% reservoir labels
\node[engred, font=\scriptsize, anchor=south] at (0,3.0) {heat to hot space $q_h$};
\node[engblue, font=\scriptsize, anchor=north] at (0,-3.0) {heat from cold space $q_c$};
\node[engorange, font=\scriptsize, anchor=south] at (4,1.0) {work $w_c$ in};
\end{tikzpicture}
\caption[Component schematic of the vapour-compression refrigerator]{The four
components of a vapour-compression refrigerator arranged as a loop. Saturated
vapour leaves the evaporator (state 1), the compressor raises its pressure
(state 2), the condenser rejects heat $q_h$ to the warm surroundings and
condenses the refrigerant (state 3), the expansion valve throttles it back to
low pressure (state 4), and the evaporator absorbs heat $q_c$ from the cold
space to close the loop. The same hardware run as a heat pump simply treats the
condenser side as the useful output.}
\label{fig:vol04ch03:vapor-compression-schematic}
\end{figure}
